% Canonical mathematical construction. Protocol applications cite these results.
\section{Coupled Hodge Encryption}
\label{sec:coupled-harmonic-construction}

\subsection{The private presentation and its message space}

The private presentation supplies the cellular differentials, inner products
and compatible integral stalk lattices. These data determine the adjoints,
Hodge operator and integral harmonic messages. We now fix a concrete family
in which one integer matrix reconstructs all of these data. Its standard
Euclidean metrics and integer stalks are part of the profile.

For a harmonic message \(h\), an exact/coexact opening gives the cochain
\begin{equation}
 b=h+\delta_0a+\delta_1^*u,
 \qquad h\in C^1_{\mathbb Z}(F)\cap\ker\Delta_{1,F}.
 \label{eq:coupled-hodge-encryption}
\end{equation}
The harmonic projector recovers \(h\). Preservation of the public homotopy
class will follow from an equivalence of the underlying cochain
interpretations; the integral lattice supplies the message space.

\subsection{Shared integral stalk constraints}

Fix \(B\in\mathbb Z^{m\times n}\). The coupled cellular family has one scalar
exact edge, \(n\) scalar information edges, \(m\) scalar parity edges, and
one scalar face stalk per row of \(B\). Its actual restriction maps give
\[
 C^1(F_B)=\mathbb R\oplus\mathbb R^n\oplus\mathbb R^m,
 \qquad
 \delta_{0,B}(a)=(a,0,0),\qquad
 \delta_{1,B}(a,x,y)=y-Bx.
\]
The information edges contribute to shared face equations, so the stalk maps
couple different edge cells. Standard integer stalk lattices are preserved by
these maps. All cochain spaces use their fixed Euclidean metrics.

These definitions reconstruct the full admitted geometry from the dimensions
and \(B\): the cellular maps, standard metric and integral stalks, harmonic
space and graph basis. Directly from the two differentials and their adjoints,
\[
 \Delta_B(a,x,y)=\bigl(a,-B^{\mathsf T}(y-Bx),y-Bx\bigr),
\]
\[
 \Delta_B=
 \begin{pmatrix}
  1&0&0\\0&B^{\mathsf T}B&-B^{\mathsf T}\\0&-B&I
 \end{pmatrix}.
\]
Each of \(\delta_{1,B}\) and \(\Delta_B\), in these fixed coordinates,
determines every entry of \(B\). Canonical \(B\) bytes can therefore bind this
entire profile without an independently chosen operator. A general Hodge
matrix does not determine arbitrary integral lattices, coefficient bases or
the separate differentials. Varying those data, the topology or the metric
requires an explicit enlarged reconstruction profile.

\begin{proposition}[Intrinsic harmonic graph lattice]
\label{prop:coupled-harmonic-graph}
The harmonic subspace and its intrinsic integral lattice are exactly
\[
 W_B=\{(0,x,Bx):x\in\mathbb R^n\},\qquad
 \Lambda_B=\{(0,z,Bz):z\in\mathbb Z^n\}.
\]
The columns of \(G_B=[0;I;B]\) present the entire lattice \(\Lambda_B\), with
Gram matrix \(G_B^{\mathsf T}G_B=I+B^{\mathsf T}B\). For target \((a,v,w)\),
the squared distance to the point with coefficient vector \(z\) is
\[
 a^2+\|v-z\|^2+\|w-Bz\|^2.
\]
The embedded harmonic subspace determines \(B\).
\end{proposition}

\begin{proof}
The harmonic equations are \(\delta_{0,B}^*(a,x,y)=0\) and
\(\delta_{1,B}(a,x,y)=0\), which give \(a=0\) and \(y=Bx\).
Integrality of the information coordinates forces \(x\in\mathbb Z^n\);
integer \(B\) then makes every parity coordinate integral. The information
coordinates also prove independence of the graph generators. Their inner
products and distances give the displayed formulas. Applying the graph
description to the information-coordinate unit vectors recovers each column
of \(B\).
\end{proof}

\begin{proposition}[The full intrinsic lattice and its coefficient charts]
\label{prop:coupled-intrinsic-lattice}
The graph map gives a real linear equivalence
\(g_B:\mathbb R^n\simeq W_B\) and an integer linear equivalence
\(e_B:\mathbb Z^n\simeq\Lambda_B\). The integer chart is the original
harmonic-lattice point encoding. Restricting the original subgroup to its
harmonic ambient space identifies it with the integer span of the real graph
basis. It is discrete, spans \(W_B\) over \(\mathbb R\), and has integral
rank \(n\).
\end{proposition}

\begin{proof}
Information coordinates invert the real graph map. The complete integral
graph characterization in Proposition~\ref{prop:coupled-harmonic-graph}
gives the integer chart and its surjectivity. Transporting the standard
coefficient bases gives the two graph bases. An element belongs to the
integer span exactly when all its graph-basis coordinates are integers,
which is the original intrinsic membership condition. The integer span of
a real basis is discrete and spans its ambient space. Its integer basis has
\(n\) elements.
\end{proof}

% Requires tikz-cd. All four maps use the original coupled geometry.
\begin{figure}[htbp]
\centering
\begin{tikzcd}[row sep=2.8em, column sep=5.0em]
 \mathbb Z^n \arrow[r, "e_B", "\simeq_{\mathbb Z}"']
     \arrow[d, "\text{scalar inclusion}"']
 & \Lambda_B \arrow[d, "\text{harmonic inclusion}"] \\
 \mathbb R^n \arrow[r, "g_B", "\simeq_{\mathbb R}"'] & W_B
\end{tikzcd}
\caption{The integer and real coefficient charts describe the same graph
\(g_B(z)=(0,z,Bz)\). The upper chart is onto the entire original intrinsic
lattice, rather than a finite-index sublattice. Inside \(W_B\), its image is
exactly the integer span of the real graph basis; its rank is \(n\) and its
inherited topology is discrete.}
\label{fig:coupled-lattice-chart}
\end{figure}


Full real span and equality with the intrinsic integral subgroup are separate
conditions in a general harmonic-lattice presentation. Both are proved here
for the same coupled graph basis; no auxiliary message lattice is chosen.
The integer native profile also derives coexact integrality from
\(\delta_{1,B}^*u=(0,-B^{\mathsf T}u,u)\). General integral-structure laws
preserve the differentials and do not assert adjoint preservation of arbitrary
chosen lattices. The finite mechanism uses this actual standard metric and
integral profile.

The planes \(W_B\) lie in the Grassmannian of \(n\)-planes in the exact-zero
space \(\mathbb R^n\oplus\mathbb R^m\). They occupy the open chart where
projection onto the information coordinates is invertible. The columns of
\([I;B]\) have information Pl\"ucker minor \(1\), and the minors obtained by
replacing one information row recover the entries of \(B\) up to sign. Thus
the full projective point determines the operative coupling in this chart.
The integer matrices select the points carrying our intrinsic graph lattices.

\subsection{Constructed homotopy-class and lattice preservation}

Helfer's homotopy equivalences compare interpretations of a declared
signature \cite[Definition~7.1]{helfer2019fol}. Here the sort is the actual
cochain complex and the closed tests concern nonzero homology. The metric and
Hodge adjoints are additional data, rather than operations whose covariance
is asserted by the semantic signature.

\begin{theorem}[Same public class and message coefficients]
\label{thm:coupled-public-class}
For integer matrices \(B,C\) of the same dimensions, define
\[
 S_{B,C}(a,x,y)=(a,x,y+(C-B)x).
\]
Together with identity maps in degrees zero and two, this is a strict
isomorphism between the actual cellular cochain complexes. It induces a
typed Helfer equivalence and equality of their closed theories in the declared
homological test fragment. Its restriction gives an additive equivalence
\(\Lambda_B\simeq\Lambda_C\) preserving the integer coefficient vector:
\[
 S_{B,C}(0,z,Bz)=(0,z,Cz).
\]
\end{theorem}

\begin{proof}
The inverse is \(S_{C,B}\), and direct substitution gives
\[
 S_{B,C}\delta_{0,B}=\delta_{0,C},\qquad
 \delta_{1,C}S_{B,C}=\delta_{1,B}.
\]
The resulting strict complex isomorphism supplies the typed semantic
equivalence. Proved invariance of the declared closed tests gives theory
equality. The integral graph description in
Proposition~\ref{prop:coupled-harmonic-graph} gives the lattice equivalence
and its coefficient formula.
\end{proof}

The shear is an isomorphism of the full cochain interpretations. It mixes
edge cells and need not define a stalkwise cellular morphism on the fixed
scalar-cell decomposition. It generally changes lengths, so it does not
transport adjoints, orthogonal projections or the noise
in~\eqref{eq:coupled-hodge-encryption} merely by being a cochain isomorphism.

For an integral chain isometry compatible with the harmonic-lattice encodings,
sample-level transport and recovery give
\[
 T_1\operatorname{Enc}_F(z,e)
   =\operatorname{Enc}_G(z,T_{\mathrm{noise}}e),\qquad
 \operatorname{Dec}_G(T_1b)=z=\operatorname{Dec}_F(b)
\]
on encryptions of \(z\). The underlying cochain equivalence preserves the
declared semantic homotopy class. Transport material may be private; this
correctness equation alone supplies no public translation algorithm.

\begin{proposition}[Native ciphertext transport]
\label{prop:coupled-native-transport}
The real coordinate map
\[
 M_B(a,z,u)=(a,z-B^{\mathsf T}u,Bz+u)
\]
is a linear equivalence. Let \(L_{B\to C}=M_CM_B^{-1}\). For every original
integral Hodge sample,
\[
 L_{B\to C}\operatorname{Enc}_B(z,a,u)
   =\operatorname{Enc}_C(z,a,u).
\]
The transports have identity, inverse and composition laws, and generated
integral samples remain generated integral samples with unchanged
\(z,a,u\).
\end{proposition}

\begin{proof}
For every real coordinate vector \(q\), the two coupling terms cancel in
\(\langle q,M_Bq\rangle=\langle q,q\rangle\). Thus \(M_B\) is injective,
hence bijective in the common finite-dimensional carrier. The original
sample's coordinate formula identifies its observed cochain with
\(M_B(a,z,u)\). The transport and composition formulas follow by cancellation
of the coordinate equivalences.
\end{proof}

\begin{proposition}[Harmonic recovery square]
\label{prop:coupled-recovery-square}
Let \(P_B\) denote harmonic projection and
\(g_B(z)=(0,z,Bz)\). For every real coordinate triple \(\xi=(a,z,u)\) and
every real cochain \(b\),
\begin{align*}
 P_B M_B\xi&=g_B(z),&
 P_C L_{B\to C}b&=S_{B,C}P_Bb,\\
 \operatorname{info}_C P_C L_{B\to C}b&=\operatorname{info}_B P_Bb.
\end{align*}
On harmonic vectors \(L_{B\to C}\) agrees with \(S_{B,C}\), including
their restrictions to the intrinsic integral lattices.
\end{proposition}

\begin{proof}
The actual assembly is
\(M_B(a,z,u)=g_B(z)+\delta_{0,B}a+\delta_{1,B}^*u\).
Harmonic projection fixes the first summand and kills the other two.
Every real \(b\) has the unique form \(M_B\xi\). Substitution of
\(L_{B\to C}M_B=M_C\) and \(S_{B,C}g_B=g_C\) gives the square.
Information coordinates of \(g_B(z)\) are \(z\), which proves the final
identity. Taking zero noise proves agreement on harmonic vectors; restricting
to integer \(z\) gives the lattice statement.
\end{proof}

% Canonical core-geometry figure. Requires tikz-cd.
% Lean correspondence: HCP.CoupledGeometry; HCP.CoupledSemantics.
\begin{figure}[htbp]
\centering
\begin{tikzcd}[row sep=2.8em, column sep=3.2em]
 & \mathcal D_{\mathbb R}
     \arrow[dl, "M_B"'] \arrow[dr, "M_C"] & \\
 C^1(F_B) \arrow[rr, "L_{B\to C}"] \arrow[d, "P_B"']
   && C^1(F_C) \arrow[d, "P_C"] \\
 W_B \arrow[rr, "S_{B,C}|_{W_B}"]
     \arrow[dr, "\operatorname{info}_B"']
   && W_C \arrow[dl, "\operatorname{info}_C"] \\
 & \mathbb R^n &
\end{tikzcd}

\medskip
\begin{tikzcd}[row sep=2.2em, column sep=3.2em]
 \mathcal I(K_B) \arrow[rr, "\simeq_{\mathrm H}"]
       \arrow[dr, "\operatorname{Th}"']
   && \mathcal I(K_C) \arrow[dl, "\operatorname{Th}"] \\
 & \alpha &
\end{tikzcd}
\caption{The core recovery diagram and its semantic class.
Here \(\mathcal D_{\mathbb R}=\mathbb R\oplus\mathbb R^n\oplus\mathbb R^m\),
\(W_B=\ker\Delta_B\), and \(\operatorname{info}_B(0,z,Bz)=z\).
The upper triangle assembles the same coordinates, the square projects to
harmonic messages, and the lower triangle reads their coefficients.
The separate semantic triangle uses the actual cochain complexes
\(K_B=C^\bullet(F_B)\), their typed interpretations, and the common declared
homological theory \(\alpha\). The semantic equivalence is induced by
\((1,S_{B,C},1)\). Its shear acts on harmonic messages in the
large diagram; native transport acts on ciphertexts.}
\label{fig:coupled-core-diagrams}
\end{figure}


This square holds without an isometry because the tailored native map
preserves the harmonic and noise summands separately. The semantic shear
cannot replace its ciphertext arrow. For \(B=0,C=1,b=(0,0,1)\),
\[
 S_{B,C}P_Bb=0,\qquad P_CS_{B,C}b=(0,1/2,1/2).
\]

An exact rational refinement computes the inverse by the
determinant/adjugate formula. Its real realization equals \(L_{B\to C}\).
This reference algorithm has no proved polynomial running-time bound.
Usable transport to the known zero presentation reveals the message
coordinates: the information component of
\(L_{B\to0}\operatorname{Enc}_B(z,a,u)=(a,z,u)\) is \(z\).

The HCP matrix in this family is the actual cellular Hodge Laplacian. Its
kernel is \(W_B\), and its operative recovery map is the original harmonic
projector. Thus the semantic interpretation, harmonic lattice, message and
observed cochain refer to the same presentation.

\subsection{Canonical bytes and direct recovery}

The versioned presentation domain \texttt{HODGE\_COUPLING\_V1} fixes this
family's scalar stalks, standard integral structures and Euclidean metric.
The bytes record both dimensions followed by the signed integer entries in
row-major order. Natural magnitudes use little-endian base-256 digits, each
tagged to distinguish it from the terminating byte. An integer adds a sign
byte. Canonical re-encoding rejects redundant digits and negative zero.
The executable decoder roundtrips the actual matrix; injectivity, rejection
of wrong dimensions or domain, trailing data and every proper truncation are
proved. Encoded magnitude length is logarithmic in its value.

For the constructed integral opening profile, \(z,a,u\) are integer-valued and
the original Hodge ciphertext has coordinates
\begin{equation}
 b=\operatorname{Enc}_B(z,a,u)=(a,\ z-B^{\mathsf T}u,\ Bz+u).
 \label{eq:coupled-integer-encryption}
\end{equation}
An executable integer vector casts to exactly this original observation. The
full-ciphertext codec uses the distinct domain
\texttt{HODGE\_CIPHERTEXT\_V1}, explicit matrix dimensions, and exact,
information and parity coordinates in that order. Its decoder reconstructs
the same native sample's full observed cochain; the existing harmonic
projector recovers its original message. Canonical roundtrip, injectivity and
rejection of malformed or mismatched byte streams are proved.

Write \(\kappa(B)\) for the private matrix bytes and \(E_B(z,a,u)\) for
the ciphertext bytes. The sender computes \eqref{eq:coupled-integer-encryption}
and encodes the resulting cochain. The recipient parses \(\kappa(B)\) and
the ciphertext in the same dimensions. For parsed coordinates \((a,x,y)\),
the recipient applies the exact rational inverse of \(M_B\). Its message
coordinate \(v\) is the unique solution of
\[
 (I+B^{\mathsf T}B)v=x+B^{\mathsf T}y.
\]
This equation follows by eliminating the coexact coordinate from the two
blocks of the ciphertext. Its matrix is positive definite. The recipient
returns \(v\) when all its coordinates are integers and rejects a parse
failure or a nonintegral result. Write \(D_B\) for recovery with a parsed
matrix and \(D_{m,n}^{\mathrm{bytes}}\) for recovery with private matrix bytes.

\begin{proposition}[Finite encryption and exact byte recovery]
\label{prop:coupled-byte-recovery}
For every integer \(B,z,a,u\),
\[
 D_{m,n}^{\mathrm{bytes}}\bigl(\kappa(B),E_B(z,a,u)\bigr)
 =D_B(E_B(z,a,u))=\operatorname{some}(z).
\]
For every byte string \(\beta\),
\[
 D_B(\beta)=\operatorname{some}(z)
 \quad\Longleftrightarrow\quad
 \exists a\in\mathbb Z,\ u\in\mathbb Z^m:\ \beta=E_B(z,a,u).
\]
On every rational cochain, the rational coefficient algorithm equals the
original harmonic-lattice partial decoder. Thus byte recovery, including
rejection, refines canonical parsing followed by the same harmonic-lattice
decoding operation.
\end{proposition}

\begin{proof}
The matrix codec returns \(B\) on \(\kappa(B)\). The rational inverse
realizes the original real inverse, whose information
coordinates are the coordinates of the projected harmonic graph. They are
integral exactly when that projection belongs to \(\Lambda_B\). For an
integer observation \((a,x,y)\) and an accepted integer message \(z\), its
coexact opening is \(u=y-Bz\), which is also integral. The original inverse
identities give \(M_B(a,z,u)=(a,x,y)\). Canonical parsing identifies exactly
the corresponding byte encoding. Conversely every such generated observation
recovers \(z\) by the same inverse and integer check.
\end{proof}

For fixed \(B\), the ciphertext also fixes both openings. Its exact
coordinate is \(a\), and \(u=y-Bz\). Two integer triples encrypted under the
same \(B\) therefore give equal bytes only when all three entries agree.
This fixed-key uniqueness concerns the ciphertext bytes directly. The
prime-field refinement below concerns a supplied assignment to the circuit
wires.

\Needspace{10\baselineskip}
\paragraph{Executing the byte round trip.}
Take \(B=(1\ 1)\), \(z=(2,-1)\), \(a=3\) and \(u=4\). The encoder produces
the canonical bytes of \((3,-2,-5,5)\). After parsing those bytes and
\(\kappa(B)\), the recipient computes the solution of
\[
 \begin{pmatrix}2&1\\1&2\end{pmatrix}v=\begin{pmatrix}3\\0\end{pmatrix},
 \qquad v=\begin{pmatrix}2\\-1\end{pmatrix}.
\]
The Lean executable returns these original coefficients and checks rejection
of a truncated private key, a truncated ciphertext and a canonically encoded
cochain with nonintegral projected coefficients. The command in
\cref{app:formal-boundaries} reproduces the execution.

\subsection{Finite-field refinement of the ciphertext equations}
\label{sec:coupled-field-refinement}

For bounded integer inputs, the same ciphertext map
\eqref{eq:coupled-integer-encryption} has a rank-one constraint
representation. The core uses \(2mn\) gates for the products
\(B_{ij}u_i\) and \(B_{ij}z_j\), then \(n+m+1\) affine equations for the
observed coordinates. The input wires \(B,z,a,u\) have absolute value at
most \(K\). Product wires have bound \(K^2\), the observed exact coordinate
has bound \(K\), and the information and parity coordinates have respective
bounds \(K+mK^2\) and \(K+nK^2\).

For a wire with field value \(x\) and bound \(K_w\), put
\(t_w=\lfloor\log_2(2K_w)\rfloor+1\) when \(K_w>0\), and \(t_w=1\)
when \(K_w=0\). Introduce two vectors of Boolean digits of length \(t_w\).
Writing their binary sums as \(S^+\) and \(S^-\), impose
\[
 x+K_w=S^+,\qquad S^++S^-=2K_w .
\]
These are two affine equations. Each Boolean digit \(e\) has the rank-one
equation \(e(e-1)=0\).

\begin{proposition}[Exact lift of the bounded ciphertext circuit]
\label{prop:coupled-field-refinement}
Every bounded integer encryption sample gives an explicit satisfying
assignment to the core and range gates over \(\mathbb F_p\). Conversely,
suppose \(p\) is prime and
\[
 p\ge 2^{t_w+1}\quad\text{for every wire }w,\qquad
 p>\max\{2K^2,\,2(K+mK^2),\,2(K+nK^2),\,2K\}.
\]
Every satisfying field assignment has a unique bounded integer lift of its
core wires. The lifted observed coordinates obey the original integer
ciphertext equations.
\end{proposition}

\begin{proof}
Use the integer sample and its actual products for the core assignment.
For every integer wire value \(v\), write \(v+K_w\) and \(K_w-v\) in
binary. Both lie in \([0,2K_w]\), so the digits fit and satisfy the
equations after reduction.

In the converse direction, each field digit is \(0\) or \(1\).
The sum of the two binary values is less than \(2^{t_w+1}\).
The congruence with \(2K_w\) is therefore an integer equality.
The difference \(S^+-K_w\), computed in the integers, is a lift of \(x\)
in \([-K_w,K_w]\). Two such lifts differ by a multiple of \(p\) and have
difference of absolute value at most \(2K_w<p\), so the lift is unique.
The core gate residuals of these lifts are divisible by \(p\).
The wire bounds give the displayed maximum as an upper bound on their
absolute values. Hence the residuals vanish, giving the integer products
and \eqref{eq:coupled-integer-encryption}.
\end{proof}

Uniqueness is relative to the full field assignment, including its private
wires. Paper 2 uses the circuit as the arithmetic component of its proof
relation and handles the public certificate and transfer checks.

\subsection{A finite private carrier and canonical relation}
\label{sec:coupled-finite-carrier}

Fix a program descriptor \(p=(m,n,K,\ell)\), with integer-coordinate bound
\(K\in\mathbb N\) and credential-salt length \(\ell\). The private carrier is
\[
 w=(S,B,z,a,u,\rho,r),\qquad
 S,B\in\mathbb Z^{m\times n},\quad z\in\mathbb Z^n,\quad
 a\in\mathbb Z,\quad u\in\mathbb Z^m.
\]
Here \(S\) is a source presentation, \(B\) is the operative target
presentation, \(\rho\) is an \(\ell\)-byte credential salt and \(r\) is a
separate 32-byte operator nonce. A bounded witness has every entry of
\(S,B,z,a,u\) of absolute value at most \(K\). Dimensions, lengths and these
bounds make its carrier finite. Both presentations use the same fixed
cellular, metric and integral profile.

Write \(\operatorname{enc}(S)\) for its canonical matrix bytes and
\[
 \operatorname{credbytes}(\rho,S)
   =\mathsf{dom}_{\mathrm{cred}}\Vert
     \operatorname{frame}(\rho)\Vert\operatorname{enc}(S).
\]
The credential domain is distinct from the matrix and ciphertext domains.
A fixed byte hash \(H\) maps byte strings to 32-byte names. The source
credential is \(H(\operatorname{credbytes}(\rho,S))\); it names fixed source
opening data rather than an identity. The private operator commitment and
three-name public relation are defined below. Their codecs retain the
original matrices and coordinates, not only a profile name.

For a sampling diagnostic, the integer-box sampler chooses the first five
integer components independently and uniformly from \(\{-K,\ldots,K\}\)
and chooses a uniform \(\ell\)-byte salt. Its native seed space has size
\[
 N=(2K+1)^{2mn+n+m+1}\,256^\ell.
\]
Each seed has mass \(N^{-1}\), and injective seed decoding gives the same
mass to its decoded native witness. Every supported witness satisfies the
selected bounds. An additional 32-byte nonce or fixed cryptographic tape
can be sampled independently. The box law includes degenerate matrices and
zero messages; its cardinality alone gives no recovery bound.

\paragraph{Scalar sampling diagnostic.}
Uniform finite coordinates alone do not give a small recovery probability.
In the original scalar sampler with bound one, a public predictor parses the
canonical ciphertext and recovers \(z\) with probability exactly \(7/9\).
Writing the scalar observation as \((a,x,y)\), the predictor returns
\(\operatorname{sign}(x)\) when \(x\ne0\), returns
\(\operatorname{sign}(y)\) when \(x=0\) and \(|y|=2\), and returns zero
otherwise. It is correct on 21 of the 27 equally likely triples
\((B,z,u)\in\{-1,0,1\}^3\). The calculation uses the full original seed law
and holds for every credential salt length, also with an independent nonce
or fixed cryptographic tape appended. It ignores hashes, the private source
presentation and the exact opening. A known zero message with \(u\ne0\)
also exposes a reused scalar operator through \(x=-Bu\), \(y=u\), hence
\(B=-x/y\). These computations concern coefficient and operator recovery.

\subsection{The fresh-sign encryption profile}
\label{sec:coupled-sampling}

\paragraph{Parameters and key generation.}
Fix a public integer matrix \(A\in\mathbb Z^{m\times n}\) with
\(A^{\mathsf T}A=dI\), \(d>0\), and a public integer exact coordinate \(a\).
Thus \(n\leq m\) and each column of \(A\) is nonzero. The payload is \(n\)
bits, represented by signs \(s\in\{-1,1\}^n\) with \(0\) mapped to \(-1\)
and \(1\) to \(1\). Sample \(n\) independent uniform secret signs \(k\) for
one encryption. The signs determine the private presentation
\(C_k=A\operatorname{diag}(k)\) injectively, and can be expanded to the
canonical matrix bytes when needed. Sender and recipient need the same
secret key. Setup for distributing it lies outside the mechanism. The
separate commitment nonce in \cref{sec:native-random-oracle} is part of the
operator commitment opening. Decryption takes the sign key and ciphertext.

\paragraph{Encryption.}
For message signs \(s\) and key \(k\), set
\[
 z=ds,\qquad u=C_ks=A(k\mathbin\odot s),
\]
where \(\odot\) is coordinatewise multiplication. Encode the full original
Hodge sample, which has the exact normal form
\[
 \operatorname{Enc}_{C_k}(ds,a,C_ks)
   =\bigl(a,0,(d+1)A(k\mathbin\odot s)\bigr),
\]
using the ciphertext codec above. Core encryption is deterministic given
the public parameters, message and sign key. It records \(1+n+m\)
integer coordinates as well as the dimension and domain framing. The first
coordinate is \(a\), the next \(n\) vanish, and the final coordinates satisfy
\(\lvert y_i\rvert\leq (d+1)\sum_j\lvert A_{ij}\rvert\). The versioned
instance adds the three 32-byte names and their framing.

\paragraph{Decryption and rejection.}
Given \(C_k\), exact private recovery returns the integer coefficients
\[
 z=\frac{x+C_k^{\mathsf T}y}{d+1}.
\]
Taking the signs of \(z=ds\) gives the original bits. Equivalently, on fresh
ciphertexts, compute the public mask and unmask it with the secret key:
\begin{equation}
 t=\frac{A^{\mathsf T}y}{d(d+1)}=k\mathbin\odot s,
 \qquad s=k\mathbin\odot t.
 \label{eq:fresh-sign-mask}
\end{equation}
The checked sign decoder parses the canonical bytes, uses the signs of
\(A^{\mathsf T}y\) as a candidate \(t\), and checks that encoding
\((a,0,(d+1)At)\) reproduces the input bytes. It returns \(k\mathbin\odot t\)
on success. A failed parse or comparison returns failure.

\begin{proposition}[Exact acceptance in the fresh-sign profile]
\label{prop:coupled-sign-checked-recovery}
For each fixed public profile with \(A^{\mathsf T}A=dI\) and \(d>0\),
each sign key \(k\), and all byte strings \(\beta\), the checked decoder
returns \(s\) if and only if \(\beta\) is exactly the canonical byte encoding
of \(\operatorname{Enc}_{C_k}(ds,a,C_ks)\).
\end{proposition}

\begin{proof}
On generated bytes the transpose calculation in \eqref{eq:fresh-sign-mask}
recovers the mask. The normal form then re-encodes to the given bytes, and
the key cancels in unmasking. Conversely, on success the parsed bytes were
equal to the normal form for the candidate mask \(t\). The returned signs
\(s=k\mathbin\odot t\) recreate that same mask and hence those exact bytes.
\end{proof}

For dense \(A\), computing each intermediate vector once gives
\(O(mn)\) scalar arithmetic operations for encryption, checked decoding or
the specialized coefficient formula in positive dimensions. These counts
treat integer operations individually and omit byte encoding and parsing.
The general recovery executable uses exact rational matrix inversion.
Both \(A=I\), \(d=1\), and the dense
matrix \(\left(\begin{smallmatrix}1&1\\1&-1\end{smallmatrix}\right)\),
\(d=2\), instantiate the profile.

\begin{proposition}[Chordal separation of the hidden planes]
\label{prop:coupled-chordal-distance}
In this fresh-sign family, let \(P_k\) be the harmonic projector onto
\(W_{C_k}\). Use chordal distance \(d_{\mathrm{ch}}(W,W')=\|P_W-P_{W'}\|_F/\sqrt 2\)
at the fixed ambient metric. For any two sign keys \(k,\ell\),
\[
 d_{\mathrm{ch}}(W_{C_k},W_{C_\ell})^2
    =\frac{4d}{(1+d)^2}\,d_H(k,\ell),
\]
where \(d_H\) counts coordinates on which the two signs differ.
\end{proposition}

\begin{proof}
With \(G_k=[0;I;C_k]\), the Gram identity gives
\(P_k=(1+d)^{-1}G_kG_k^{\mathsf T}\). The exact block vanishes, and the
information and parity diagonal blocks of \(G_kG_k^{\mathsf T}\) are
\(I\) and \(AA^{\mathsf T}\). Hence only the off-diagonal blocks vary.
Writing \(D_k=\operatorname{diag}(k)\), the squared Frobenius norm of
\(C_k-C_\ell=A(D_k-D_\ell)\) is \(4d\,d_H(k,\ell)\). There are two
transposed copies of this block in the projector difference. Dividing their
squared norms by \(2(1+d)^2\) gives the identity.
\end{proof}

All of these planes share the public homotopy class. Their exact distances
describe the geometry of the ambiguity, and in \cref{prop:coupled-cvp-candidate-count}
the same planes give distinct closest-vector answers to one ciphertext.

\paragraph{Executing the fresh-sign profile.}
Use this dense matrix, Boolean message \((1,0)\) and key signs \((-1,1)\).
With \(a=3\), the operative matrix and openings give
\[
 C_k=\begin{pmatrix}-1&1\\-1&-1\end{pmatrix},\qquad
 z=(2,-2),\qquad u=(-2,0),\qquad b=(3,0,0,-6,0).
\]
The encoder emits the bytes of \(b\). The recipient reads the canonical
private matrix and ciphertext bytes and recovers \((2,-2)\). Testing which
coefficients are positive produces \((1,0)\). For every \(d>0\), the same
decoder returns the original bits on generated ciphertext bytes by
\cref{prop:coupled-byte-recovery} and the sign rule for \(ds\). The executable
also runs the checked sign decoder. It returns the same bits and rejects
the canonical ciphertext with exact coordinate changed to \(4\) at the
declared profile \(a=3\). The probability laws below use fresh keys.

\subsection{Public ciphertext, integrity and key reuse}
\label{sec:coupled-sign-reuse}

The fresh-sign ciphertext discloses \(a\), \(u=y/(d+1)\) and the mask \(t\)
in \eqref{eq:fresh-sign-mask}. Given the public profile, the mask determines
the entire ciphertext. The ciphertext-only mechanism is therefore a
one-time pad in a reversible Hodge encoding. For any fixed message,
multiplication by its sign vector permutes the uniform key space. Thus the
full canonical ciphertext bytes have the same distribution for every fixed message, as does
any deterministic digest of those bytes. With independent fresh sign keys,
the ciphertext-only law for a fixed sequence of messages is the product of
these laws.

For an independent uniform Boolean message and fresh uniform key, the actual
joint law factors into a uniform message independent of ciphertext bytes.
Every total complete-message decoder succeeds with probability exactly
\(2^{-n}\), including a decoder using a separately sampled fixed random tape
or a deterministic ciphertext digest. Positive \(d\) identifies correct
encoded coefficient recovery with bit-vector recovery. An arbitrary partial
integer-coefficient extractor, allowing invalid outputs or failure, succeeds
with probability at most \(2^{-n}\), also with the independent tape.

Every sign key \(\ell\) decrypts a given valid ciphertext to the unique
message \(\ell\mathbin\odot t\). For \(\ell\ne k\), this is a different
message whose encryption under \(\ell\) produces the same bytes. Thus
successful parsing and checked decoding provide no authentication. The
core ciphertext is malleable. Anyone can replace
\(t\) with \(r\mathbin\odot t\) in the normal form for a chosen sign
vector \(r\), producing a valid ciphertext of \(r\mathbin\odot s\). The
message bits at the negative entries of \(r\) are flipped.

\begin{proposition}[Exact consequences of sign-key reuse]
\label{prop:coupled-sign-reuse}
For two fresh-sign ciphertexts with the same public \(A,d\) and the same
secret sign key \(k\), let \(t_1,t_2\) be their public masks and \(s_1,s_2\)
their respective messages. Their exact coordinates \(a_1,a_2\) may differ.
The ciphertexts determine \(t_1\mathbin\odot t_2=s_1\mathbin\odot s_2\),
which gives the bitwise XOR by mapping equal signs to zero. Exactly \(2^n\)
ordered message pairs are compatible with the two byte strings under a
common sign key. One known message recovers that key and every other message
encrypted under it in the same profile.
\end{proposition}

\begin{proof}
Apply \eqref{eq:fresh-sign-mask} to both byte strings and multiply. The key
signs cancel since their squares are one. Conversely, for each candidate
\(s_1'\), the key \(t_1\mathbin\odot s_1'\) and the second message
\(s_1'\mathbin\odot t_1\mathbin\odot t_2\) give the two observed masks.
Their canonical encryptions reproduce both original byte strings. These
are all the compatible pairs and have distinct first messages. If \(s_1\)
is known, the recovered key is \(t_1\mathbin\odot s_1\), and substitution
in \eqref{eq:fresh-sign-mask} gives every other message.
\end{proof}

These computations take the two ciphertext byte strings alone. Changing
the operator commitment nonce leaves the equations intact. A known or
chosen message under a reused key therefore compromises other ciphertexts
under that key. Two independently drawn uniform keys can coincide with
probability \(2^{-n}\); the independent ciphertext law above averages over
those draws.

For two chosen candidate messages \(s_0,s_1\), compute
\(k_i=t\mathbin\odot s_i\) and \(C_i=A\operatorname{diag}(k_i)\).
A published raw operator digest \(H(\operatorname{enc}(C_k))\) permits testing both
candidates with two forward hashes. The actual byte parser and public
transpose/parity arithmetic recover the masked signs and construct each
candidate's canonical operator bytes. The true plaintext passes; a wrong
candidate that passes yields a collision between distinct canonical matrix
bytes. This motivates the private-nonce commitment below.

The released view studied below consists only of the versioned instance and
its full ciphertext bytes. It excludes exposed projectors, kernel or
Pl\"ucker coordinates, operator nullifiers and usable transport to a known
presentation. The exact opening \(a\) is fixed independently of the message
and fresh key. No proof or storage transcript is included in this view.

Fresh sampling is distinct from deterministic migration. Transport retains
\(z,a,u\), whereas a fresh target key prescribes \(u'=C_{k'}s\). At that
target, linearity gives the exact correction
\[
 \operatorname{Enc}_{C_{k'}}(z,a,u)
 +\operatorname{Enc}_{C_{k'}}(0,0,u'-u)
 =\operatorname{Enc}_{C_{k'}}(z,a,u').
\]
Canonical byte parsing and addition implement this same coexact correction.
The identity does not itself sample a fresh key or supply a public refresh
key.

\subsection{Private operator commitments and the versioned instance}
\label{sec:private-operator-instance}

The private operator opening consists of the canonical coupling packet and
the separate 32-byte nonce \(r\). In the fixed scalar Euclidean profile,
this packet determines the cellular maps, Hodge operator, integral structure
and intrinsic harmonic lattice. Define
\[
 \operatorname{Com}(B;r)=
 H\bigl(\mathsf{dom}_{\mathrm{op}}\Vert
   \operatorname{frame}(\operatorname{enc}(B))\Vert
   \operatorname{frame}(r)\bigr).
\]
The nonce is distinct from the source credential salt. Canonical decoding
enforces the operator domain, original matrix codec, fixed nonce width and
full consumption. Distinct accepted matrices with equal commitments give
an actual hash collision between these complete preimages, allowing their
nonces to differ.

For the private carrier \(w=(S,B,z,a,u,\rho,r)\) and fixed descriptor
\(p=(m,n,K,\ell)\), define its public instance by the three names
\begin{equation}
 x_w=\bigl(H(\operatorname{credbytes}(\rho,S)),\,
          \operatorname{Com}(B;r),\,H(\beta)\bigr),
 \qquad \beta=E_B(z,a,u).
 \label{eq:coupled-committed-instance}
\end{equation}
Only the instance and ciphertext bytes are released; neither matrix nor
opening is serialized in the public instance. The common declared semantic
class is fixed for these dimensions. The versioned private-witness,
program and instance formats use the respective domains
\texttt{HODGE\_NATIVE\_WITNESS\_V2},
\texttt{HODGE\_NATIVE\_PROGRAM\_V2} and
\texttt{HODGE\_NATIVE\_INSTANCE\_V2}. They retain the canonical matrix,
integer and salt encodings and reject the earlier raw-operator-hash formats.
The finite Boolean relation is
\[
\begin{aligned}
 R_{p,H}(x,\eta)=\mathsf{true}
 \quad\Longleftrightarrow\quad \exists w:\;&
 \eta=\operatorname{encodeWitness}(w),\\
 &w\text{ satisfies the bounds for }p,\quad x=x_w.
\end{aligned}
\]
Canonical program decoding fixes \(m,n,K,\ell\), and the checker recomputes
all three names from the decoded witness. This gives both honest-witness
acceptance and reflection of acceptance into the stated relation.

\begin{proposition}[Committed-instance recovery or a hash collision]
\label{prop:coupled-committed-recovery}
Suppose an accepted witness \(w\) names an original target \(B_0\), with
operator nonce \(r_0\), and the original ciphertext
\(E_{B_0}(z_0,a_0,u_0)\). Then its message is \(z_0\), or there are distinct
canonical operator preimages with equal hashes, or distinct canonical
ciphertext byte strings with equal hashes.
\end{proposition}

\begin{proof}
The second name gives equality of the two target commitments. Unless their
complete preimages exhibit a hash collision, canonical injectivity gives
\(B=B_0\). The third name similarly gives equality of the ciphertext bytes
unless they exhibit a byte-hash collision. In the remaining case,
Proposition~\ref{prop:coupled-byte-recovery} applied with this same private
matrix recovers both messages from the same bytes, so \(z=z_0\).
\end{proof}

For the signed family, let \(\varepsilon(0)=-1\) and \(\varepsilon(1)=1\), applied
coordinatewise to Boolean vectors, and put
\[
 B_k=A\operatorname{diag}(\varepsilon(k)),\qquad
 z=d\varepsilon(\mu),\qquad u=B_k\varepsilon(\mu).
\]
Sample the Boolean message \(\mu\in\{0,1\}^n\), key \(k\in\{0,1\}^n\)
and nonce independently and uniformly. Fix the source opening \((\rho,S)\)
and exact opening \(a\) independently of these samples. Its credential is
therefore fixed in the following fixed-hash experiment. The released view
is the pair
\[
 V=\bigl(\operatorname{encodeInstance}(x_w),E_{B_k}(z,a,u)\bigr).
\]
If the bounded checker is used, select \(K\) large enough for this finite
family and the fixed source opening.

Fix a reference matrix \(B_*\) of the same dimensions. For a deterministic
complete-message decoder \(D\), the chosen-operator test \(T_D\) samples
its own \(\mu,k\), receives either \(\operatorname{Com}(B_k;r)\) or
\(\operatorname{Com}(B_*;r)\), constructs the same ciphertext and public
view, runs \(D\), and tests whether the result is \(\mu\). Let \(p_1,p_0\)
be its real and reference acceptance probabilities, respectively, and put
\(\operatorname{Adv}_{\mathrm{op}}(T_D)=|p_1-p_0|\).

\begin{theorem}[Public-instance recovery reduces to operator hiding]
\label{thm:coupled-instance-recovery}
For \(A^{\mathsf T}A=dI\) and every such decoder \(D\),
\[
 \Pr[D(V)=\mu]\leq 2^{-n}+\operatorname{Adv}_{\mathrm{op}}(T_D).
\]
For positive \(d\), a partial integer-coefficient extractor has the same
bound using its threshold-totalized bit decoder in \(T_D\).
\end{theorem}

\begin{proof}
In the reference world, the operator root depends only on the independent
nonce and fixed matrix. The ciphertext law, with that nonce as an
independent tape, gives complete-message recovery probability exactly
\(2^{-n}\). The real world's public names and ciphertext are exactly those
computed from the sampled private witness. The two recovery events are the
acceptance events of \(T_D\), so their difference is bounded by its
advantage. For a partial coefficient extractor, replace failure by a fixed
bit vector and an integer output by its coordinatewise positive-sign bits.
When \(d>0\), every correct coefficient output gives the correct Boolean
message. Its success event is therefore contained in the totalized decoder's
success event.
\end{proof}

This is an algebraic transfer for a fixed byte hash, with the advantage of
an explicit chosen-operator test. The next experiment instead samples one
shared random oracle and bounds the actual recovery event directly.

\subsection{One-use privacy in the shared byte random oracle}
\label{sec:native-random-oracle}

The public oracle accepts arbitrary byte strings and returns independent
uniform 32-byte digests on previously unseen inputs, caching each answer.
The experiment starts with an empty cache. Before releasing any view, the
constructor samples independent uniform Boolean vectors
\(\mu,k\in\{0,1\}^n\) and an independent uniform 32-byte nonce \(r\).
Here \(\varepsilon(0)=-1\), \(\varepsilon(1)=1\) coordinatewise, and
\[
 B_k=A\operatorname{diag}(\varepsilon(k)),\qquad
 z=d\varepsilon(\mu),\qquad u=B_k\varepsilon(\mu).
\]
The source opening \((\rho,S)\), exact opening \(a\), and public matrix \(A\)
are fixed independently of these samples. Choose a common descriptor with
\(\ell=|\rho|\) and a coordinate bound \(K\) satisfying
\[
\begin{gathered}
 |S_{ij}|\leq K,\quad |A_{ij}|\leq K,\quad |d|\leq K,\quad |a|\leq K,\\
 \sum_j |A_{ij}|\leq K\qquad\text{for every row }i.
\end{gathered}
\]
Every generated private carrier then satisfies the same bounded relation:
the target entries have magnitude \(|A_{ij}|\), the message entries have
magnitude \(|d|\), and the coexact entries are bounded by the row sums.

The constructor first obtains the operator root by querying this same byte
oracle at the complete canonical preimage for
\(\operatorname{Com}(B_k;r)\). It then queries the source-credential bytes
and full ciphertext bytes, and releases
\(\operatorname{encodeInstance}(x_w)\) together with that ciphertext.
The decoder receives only this view, uniform random sampling and the public
byte oracle. It has no private-derivation interface or earlier oracle
transcript. Its query budget \(Q\) counts all of its postrelease byte-oracle
queries, on every input view and every execution path.

The proof uses the hidden-seed simulation theorem of the pinned VCVio
library \cite{vcvio2026library}. The real handler shares one cache between
public queries and the constructor's private operator derivation; the ideal
handler uses an independent derivation table. Injectivity of the canonical
packet/nonce encoding implies that any fixed byte address is an operator
opening for at most one nonce seed. A uniform 32-byte seed therefore hits
that address with probability at most \(256^{-32}\). The adaptive event
comparison charges this term once per public byte-oracle query.

\begin{theorem}[Native recovery in the random-oracle model]
\label{thm:coupled-native-rom}
Fix \(A\in\mathbb Z^{m\times n}\) with \(A^{\mathsf T}A=dI\), a fixed
source credential opening and exact opening \(a\). In the fresh-cache
experiment above, every decoder making at most \(Q\) public byte-oracle
queries satisfies
\[
 \Pr_{\mathrm{RO}}[\text{recover the complete Boolean message}]
 \leq 2^{-n}+\frac{Q+2}{256^{32}}.
\]
If \(d>0\), the same upper bound holds for a partial extractor of the
original integer harmonic coefficients \(d\varepsilon(\mu)\). The two additional
queries are the constructor's source-credential and full-ciphertext digest
calls.
\end{theorem}

\begin{proof}
The constructor makes exactly one private derivation, before any public
calls, and then two public digest calls. Uniform draws are uncharged. The
lifted decoder has no private-derivation access, so the complete program
has public-query budget \(Q+2\). Hidden-seed event comparison bounds real
recovery by ideal recovery plus \((Q+2)256^{-32}\).

In the ideal game, the first derivation returns an independent uniform
operator root. All later calls use only the public byte oracle. Projecting
away the separate derivation table leaves that same public-oracle execution.
Reindex the fresh key by the fixed-message mask
\(\varepsilon(k)\mathbin\odot\varepsilon(\mu)\). This is a permutation of the uniform
key space, and the native normal form makes the ciphertext depend only on
this mask. The ideal operator root, ciphertext, two public digests and
entire subsequent adaptive decoder execution are consequently independent
of the uniform message. Complete-message guessing has mass exactly
\(2^{-n}\). For a partial coefficient extractor, the same factorization and
positive scaling give success probability at most \(2^{-n}\); failure and
invalid integer outputs cannot increase that event.
\end{proof}

\begin{theorem}[One-use chosen-message indistinguishability]
\label{thm:coupled-native-rom-challenge}
Fix any pair of Boolean messages before the experiment. Replace the uniform
message draw above by a fair challenge bit \(b\) and encrypt the \(b\)-th
message. The key, nonce, initial cache and released view are otherwise as
above. A public adversary making at most \(Q\) adaptive byte-oracle queries
after release has
\[
 \left|\Pr_{\mathrm{RO}}[\widehat b=b]-\tfrac12\right|
 \leq \frac{Q+2}{256^{32}}.
\]
\end{theorem}

\begin{proof}
In the ideal experiment the operator root is independent and the uniform
key can be reindexed by the mask for either chosen message. The complete
public execution, including the adaptive query answers, therefore has the
same law for \(b=0\) and \(b=1\). A Boolean guess is correct with probability
\(\tfrac12\). Apply the same hidden-seed event comparison to a correct
challenge-bit guess and then to its complementary event. The query budget
remains \(Q+2\).
\end{proof}

\begin{corollary}[Native recovery with deterministic key disclosure]
\label{cor:coupled-native-rom-disclosure}
In the experiment of \cref{thm:coupled-native-rom}, also give the decoder
\(L(k)\) for a fixed deterministic function \(L\) of the sampled sign key.
The disclosure has no nonce input and makes no oracle queries. Let
\(M=|\operatorname{im} L|\). A decoder making at most \(Q\) public byte-oracle
queries on every execution path has
\[
 \Pr_{\mathrm{RO}}[\text{recover the target}]
 \leq \frac{M}{2^n}+\frac{Q+2}{256^{32}}
\]
for any injective encoding of the complete message. This includes the
Boolean message and, for \(d>0\), the original integer harmonic coefficients
with a possible failure output. If \(L\) discloses \(s\) bits, the first term
is at most \(2^{s-n}\).
\end{corollary}

\begin{proof}
The query budget remains \(Q+2\), including the two honest digest calls.
The same hidden-seed comparison adds \((Q+2)256^{-32}\) to the ideal
recovery bound. In the ideal game the mask is uniform independently of the
key, since it permutes the uniform message for each key. At a fixed mask
and disclosed value, the decoder has one output distribution, including
its adaptive public queries. The compatible keys in that disclosure fiber
give distinct targets, so the probabilities assigned to those targets
sum to at most one. Each key has weight \(2^{-n}\) at this mask, hence the
fiber contributes at most \(2^{-n}\). Sum over the \(M\) disclosed values
and average over the mask.
\end{proof}

The bounds average over one shared random oracle in the one-use
instance/ciphertext experiment. The opening nonce supplies the
\(256^{32}=2^{256}\) address-guessing denominator. The payload length \(n\)
sets the uniform-message guessing baseline. Concrete hash instantiation,
prechallenge queries, repeated oracle challenges, and additional public
proof or storage transcripts require corresponding security experiments.
The key-reuse result above already applies to the commitment-bearing view
because the ciphertext bytes are included in it.

